\question[5]
Das Drehkreuz am Eingang eines Schwimmbads wird durch den folgenden Mealy-Automaten gesteuert.
Eingaben: $m$ (Münze eingeworfen), $d$ (Person drückt gegen das Drehkreuz).

\begin{center}
\includegraphics[height=2.8cm]{\pfad/mealy_03/mealy_03.png}
\end{center}

\begin{parts}
\part Gib die Zustandsübergangstabelle an.
\part Bestimme die Ausgaben und die Zustände für die Eingabefolge $m, d, d, m, m, d$.
\part Erläutere, wie sich das Drehkreuz verhält, wenn zwei Personen hintereinander mit nur einer Münze durchgehen wollen.
\end{parts}

\begin{solutionbox}{7cm}
\begin{parts}
\part
\begin{tabular}[t]{r|cc}
 & m & d \\ \hline
$\rightarrow$ verriegelt & offen / entriegeln & verriegelt / Alarm \\
offen & offen / Münze zurück & verriegelt / verriegeln \\
\end{tabular}
\part
\begin{tabular}[t]{lll}
Eingabe & Ausgabe & neuer Zustand \\ \hline
m & entriegeln & offen \\
d & verriegeln & verriegelt \\
d & Alarm & verriegelt \\
m & entriegeln & offen \\
m & Münze zurück & offen \\
d & verriegeln & verriegelt \\
\end{tabular}
\part Nach der Münze ist das Drehkreuz offen. Die erste Person geht durch, danach ist es wieder verriegelt.
Drückt die zweite Person, bleibt es verriegelt und gibt Alarm.
\end{parts}
\end{solutionbox}
